\subsubsection{Dehomogenization and local stress recovery}
\label{sec:sc_dehomogenization}

The accuracy of the proposed hybrid formulation is further assessed by
recovering the local three-dimensional stress field in the reference
simple-cubic unit cell. The geometry and material properties are the same as
those used in Table~\ref{tab:sc_reference_results}. In particular, the solid
junction extends from the unit-cell center to $x=0.015$, whereas the remaining
portion of the member along the positive $x$ direction is represented by a
one-dimensional beam extending from $x=0.015$ to $x=0.1$. The constituent is
isotropic, with Young's modulus $70\times10^9$ and Poisson's ratio $0.3$.

A general macroscopic strain state containing all six strain components is
prescribed so that the complete local stress vector can be examined in a
single analysis. The applied macroscopic displacement, deformation matrix,
and engineering strain vector are
\begin{align}
\overline{\boldsymbol{u}}
&=
\begin{bmatrix}
-1.1056236823\times10^{-3} &
-1.6326786740\times10^{-3} &
 5.3131661811\times10^{-4}
\end{bmatrix}^{\mathsf T},
\\
\overline{\boldsymbol{C}}
&=\boldsymbol{I},
\\
\overline{\boldsymbol{\varepsilon}}
&=
\begin{bmatrix}
-1.0963464497\times10^{-4} &
 1.7514917499\times10^{-4} &
 1.9654751819\times10^{-4} &
 1.1855396374\times10^{-4} &
-1.8718415435\times10^{-4} &
-6.5452337092\times10^{-5}
\end{bmatrix}^{\mathsf T},
\end{align}
where
$\overline{\boldsymbol{\varepsilon}}
=[\varepsilon_{11},\varepsilon_{22},\varepsilon_{33},
2\varepsilon_{23},2\varepsilon_{13},2\varepsilon_{12}]^{\mathsf T}$.
The prescribed translation sets the reference position of the recovered
displacement field and does not affect the strain or stress, while
$\overline{\boldsymbol{C}}=\boldsymbol{I}$ introduces no macroscopic rotation.
The same macroscopic state is applied to the complete three-dimensional
SwiftComp model and to the Euler--Bernoulli and Timoshenko hybrid models.

The local response is compared along the centerline
\begin{equation}
\Gamma_x
=
\left\{(x,0,0)\,\middle|\,0\leq x\leq0.1\right\}.
\end{equation}
This line passes from the center of the solid junction into the member aligned
with the positive $x$ direction. The interval $0\leq x\leq0.015$ therefore
belongs to the three-dimensional junction, and $0.015<x\leq0.1$ belongs to the
beam. The vertical dashed line in
Figures~\ref{fig:sc_dehomo_S11}--\ref{fig:sc_dehomo_S12} marks this interface.

For the beam portion, OpenSG first recovers the centerline displacement,
rotation, generalized strain, and force and moment resultants from the
homogenization solution. These quantities are passed to VABS, which recovers
the three-dimensional displacement, strain, and stress fields over the square
cross-section. Eleven equally spaced axial stations are used on each member.
The stress at the cross-sectional centroid is then evaluated at each station
and interpolated to the $x$ coordinates of the SwiftComp reference data.

The junction stress is recovered on a structured C3D20 mesh containing 2,585
nodes and 448 elements. Four elements are used across the $0.01\times0.01$
cross-section, giving an element size of $0.0025$. Because the comparison line
must pass through actual finite-element nodes, the quadratic edge nodes give
13 junction nodes on $\Gamma_x$, located at intervals of $0.00125$ from
$x=0$ to $x=0.015$. The interface generalized displacements obtained from the
hybrid unit-cell solution are imposed on the six junction connection surfaces.
After solving the junction problem, stresses contributed by the adjoining
C3D20 elements are volume averaged at each centerline node. Thus, the junction
curves in the following figures are based on nodal values located exactly on
the comparison line rather than values sampled from neighboring points.

The recovered stress vector is reported in the order
\begin{equation}
\boldsymbol{\sigma}
=
\begin{bmatrix}
S_{11} & S_{22} & S_{33} & S_{23} & S_{13} & S_{12}
\end{bmatrix}^{\mathsf T}.
\end{equation}
Figures~\ref{fig:sc_dehomo_S11}--\ref{fig:sc_dehomo_S12} compare each stress
component separately. SwiftComp nodal values are shown as symbols, while the
OpenSG Euler--Bernoulli and Timoshenko results are shown as lines.

\begin{figure}[H]
\centering
\includegraphics[width=0.82\textwidth]{simple_cubic_plus_x_S11.png}
\caption{Comparison of the recovered normal stress $S_{11}$ along
$\Gamma_x$. The dashed vertical line identifies the junction--beam interface
at $x=0.015$.}
\label{fig:sc_dehomo_S11}
\end{figure}

\begin{figure}[H]
\centering
\includegraphics[width=0.82\textwidth]{simple_cubic_plus_x_S22.png}
\caption{Comparison of the recovered normal stress $S_{22}$ along
$\Gamma_x$.}
\label{fig:sc_dehomo_S22}
\end{figure}

\begin{figure}[H]
\centering
\includegraphics[width=0.82\textwidth]{simple_cubic_plus_x_S33.png}
\caption{Comparison of the recovered normal stress $S_{33}$ along
$\Gamma_x$.}
\label{fig:sc_dehomo_S33}
\end{figure}

\begin{figure}[H]
\centering
\includegraphics[width=0.82\textwidth]{simple_cubic_plus_x_S23.png}
\caption{Comparison of the recovered shear stress $S_{23}$ along
$\Gamma_x$.}
\label{fig:sc_dehomo_S23}
\end{figure}

\begin{figure}[H]
\centering
\includegraphics[width=0.82\textwidth]{simple_cubic_plus_x_S13.png}
\caption{Comparison of the recovered shear stress $S_{13}$ along
$\Gamma_x$.}
\label{fig:sc_dehomo_S13}
\end{figure}

\begin{figure}[H]
\centering
\includegraphics[width=0.82\textwidth]{simple_cubic_plus_x_S12.png}
\caption{Comparison of the recovered shear stress $S_{12}$ along
$\Gamma_x$.}
\label{fig:sc_dehomo_S12}
\end{figure}

All six stress components exhibit the same spatial trends as the SwiftComp
solution. The strongest variation occurs inside the solid junction, where the
three intersecting members produce a fully three-dimensional stress state. In
the beam region, $S_{11}$ approaches an almost uniform compressive value,
whereas the transverse normal stresses and $S_{23}$ decay toward zero. The
$S_{13}$ and $S_{12}$ components approach the smaller nonzero values associated
with the recovered transverse force resultants.

The Euler--Bernoulli and Timoshenko curves are nearly coincident for this load
state. The largest absolute OpenSG--SwiftComp difference among the beam-line
values is $0.0564$~MPa, and the largest difference among the junction-line
values is $0.3155$~MPa. The close agreement in both regions shows that the
hybrid solution transfers the macroscopic deformation consistently to the
one-dimensional members and the three-dimensional junction, while the VABS
localization reproduces the stress state in the beam region and captures the
transition from the junction response.
